% Options for packages loaded elsewhere
% Options for packages loaded elsewhere
\PassOptionsToPackage{unicode}{hyperref}
\PassOptionsToPackage{hyphens}{url}
\PassOptionsToPackage{dvipsnames,svgnames,x11names}{xcolor}
%
\documentclass[
  10pt,
  letterpaper,
  DIV=11,
  numbers=noendperiod]{scrartcl}
\usepackage{xcolor}
\usepackage{amsmath,amssymb}
\setcounter{secnumdepth}{5}
\usepackage{iftex}
\ifPDFTeX
  \usepackage[T1]{fontenc}
  \usepackage[utf8]{inputenc}
  \usepackage{textcomp} % provide euro and other symbols
\else % if luatex or xetex
  \usepackage{unicode-math} % this also loads fontspec
  \defaultfontfeatures{Scale=MatchLowercase}
  \defaultfontfeatures[\rmfamily]{Ligatures=TeX,Scale=1}
\fi
\usepackage{lmodern}
\ifPDFTeX\else
  % xetex/luatex font selection
\fi
% Use upquote if available, for straight quotes in verbatim environments
\IfFileExists{upquote.sty}{\usepackage{upquote}}{}
\IfFileExists{microtype.sty}{% use microtype if available
  \usepackage[]{microtype}
  \UseMicrotypeSet[protrusion]{basicmath} % disable protrusion for tt fonts
}{}
% Make \paragraph and \subparagraph free-standing
\makeatletter
\ifx\paragraph\undefined\else
  \let\oldparagraph\paragraph
  \renewcommand{\paragraph}{
    \@ifstar
      \xxxParagraphStar
      \xxxParagraphNoStar
  }
  \newcommand{\xxxParagraphStar}[1]{\oldparagraph*{#1}\mbox{}}
  \newcommand{\xxxParagraphNoStar}[1]{\oldparagraph{#1}\mbox{}}
\fi
\ifx\subparagraph\undefined\else
  \let\oldsubparagraph\subparagraph
  \renewcommand{\subparagraph}{
    \@ifstar
      \xxxSubParagraphStar
      \xxxSubParagraphNoStar
  }
  \newcommand{\xxxSubParagraphStar}[1]{\oldsubparagraph*{#1}\mbox{}}
  \newcommand{\xxxSubParagraphNoStar}[1]{\oldsubparagraph{#1}\mbox{}}
\fi
\makeatother


\usepackage{longtable,booktabs,array}
\usepackage{calc} % for calculating minipage widths
% Correct order of tables after \paragraph or \subparagraph
\usepackage{etoolbox}
\makeatletter
\patchcmd\longtable{\par}{\if@noskipsec\mbox{}\fi\par}{}{}
\makeatother
% Allow footnotes in longtable head/foot
\IfFileExists{footnotehyper.sty}{\usepackage{footnotehyper}}{\usepackage{footnote}}
\makesavenoteenv{longtable}
\usepackage{graphicx}
\makeatletter
\newsavebox\pandoc@box
\newcommand*\pandocbounded[1]{% scales image to fit in text height/width
  \sbox\pandoc@box{#1}%
  \Gscale@div\@tempa{\textheight}{\dimexpr\ht\pandoc@box+\dp\pandoc@box\relax}%
  \Gscale@div\@tempb{\linewidth}{\wd\pandoc@box}%
  \ifdim\@tempb\p@<\@tempa\p@\let\@tempa\@tempb\fi% select the smaller of both
  \ifdim\@tempa\p@<\p@\scalebox{\@tempa}{\usebox\pandoc@box}%
  \else\usebox{\pandoc@box}%
  \fi%
}
% Set default figure placement to htbp
\def\fps@figure{htbp}
\makeatother





\setlength{\emergencystretch}{3em} % prevent overfull lines

\providecommand{\tightlist}{%
  \setlength{\itemsep}{0pt}\setlength{\parskip}{0pt}}



 


\usepackage{xcolor, soul}
\definecolor{Light}{gray}{.90}
\sethlcolor{Light}
\let\OldTexttt\texttt
\renewcommand{\texttt}[1]{\OldTexttt{\hl{#1}}}
\usepackage{caption}
\captionsetup{font=small, labelfont=bf}
\setlength{\parindent}{0pt}
\setlength{\parskip}{1.25ex}
\KOMAoption{captions}{tableheading}
\makeatletter
\@ifpackageloaded{caption}{}{\usepackage{caption}}
\AtBeginDocument{%
\ifdefined\contentsname
  \renewcommand*\contentsname{Table of contents}
\else
  \newcommand\contentsname{Table of contents}
\fi
\ifdefined\listfigurename
  \renewcommand*\listfigurename{List of Figures}
\else
  \newcommand\listfigurename{List of Figures}
\fi
\ifdefined\listtablename
  \renewcommand*\listtablename{List of Tables}
\else
  \newcommand\listtablename{List of Tables}
\fi
\ifdefined\figurename
  \renewcommand*\figurename{Figure}
\else
  \newcommand\figurename{Figure}
\fi
\ifdefined\tablename
  \renewcommand*\tablename{Table}
\else
  \newcommand\tablename{Table}
\fi
}
\@ifpackageloaded{float}{}{\usepackage{float}}
\floatstyle{ruled}
\@ifundefined{c@chapter}{\newfloat{codelisting}{h}{lop}}{\newfloat{codelisting}{h}{lop}[chapter]}
\floatname{codelisting}{Listing}
\newcommand*\listoflistings{\listof{codelisting}{List of Listings}}
\makeatother
\makeatletter
\makeatother
\makeatletter
\@ifpackageloaded{caption}{}{\usepackage{caption}}
\@ifpackageloaded{subcaption}{}{\usepackage{subcaption}}
\makeatother
\usepackage{bookmark}
\IfFileExists{xurl.sty}{\usepackage{xurl}}{} % add URL line breaks if available
\urlstyle{same}
\hypersetup{
  pdftitle={Regression approach to representation},
  pdfauthor={Al Xin},
  colorlinks=true,
  linkcolor={blue},
  filecolor={Maroon},
  citecolor={Blue},
  urlcolor={Blue},
  pdfcreator={LaTeX via pandoc}}


\title{Regression approach to representation}
\author{Al Xin}
\date{Jun 25, 2026}
\begin{document}
\maketitle


\section{Data}\label{data}

\subsection{Loading}\label{loading}

See \texttt{20260521.qmd} for reference. Like before, the stimulus
labels are identical across all files, though the block order may
differ. For this data, subject info is aggregated, and I cannot
confidently identify which subject each recording site belongs to. Also,
like before, I removed all \texttt{NA} sites and converted the binned
MUAs to \texttt{qs2} files for faster loading.

\subsection{Predictors}\label{predictors}

As review, the stimulus labels are:

\begin{itemize}
\tightlist
\item
  \texttt{coherence}: \texttt{0} when blank, \texttt{1} when noisy,
  \texttt{2} when naturalistic.
\item
  \texttt{family}: \texttt{0} when blank, two or three digit integer
  code otherwise.
\item
  \texttt{id}: \texttt{0} when blank, integer from \texttt{1} to
  \texttt{15} otherwise.
\item
  \texttt{id\_blank}: integer from \texttt{1} to \texttt{100} if blank,
  \texttt{0} otherwise.
\item
  \texttt{repetition}: column constructed from counting the repetition
  of \texttt{coherence}, \texttt{family}, and \texttt{id} or
  \texttt{id\_blank} combo.
\end{itemize}

For regression, unique values of \texttt{id}, \texttt{id\_blank}, or
\texttt{repetition} can largely be ignored. These can basically be
considered repetitions of the unique \texttt{coherence} and
\texttt{family} combinations. There are 36 unique family values,
including the 35 texture families and a blank family.

Besides the above stimulus labels, the \texttt{block\_order}, which
takes values 1 through 8, might also be informative. If needed,
indicators like \texttt{is\_first} or \texttt{is\_last} might be a
working simplification for \texttt{block\_order}.

So, the current set of predictors are:

\begin{itemize}
\tightlist
\item
  \texttt{region}: categorical, with values \texttt{"V1"},
  \texttt{"V2"}, \texttt{"V4"}, and \texttt{"IT"}.
\item
  \texttt{coherence}: categorical, represented as dummy variables
  \texttt{is\_noisy} and \texttt{is\_natur}, with \texttt{is\_blank}
  implicit.
\item
  \texttt{family}: categorical, represented as 35 dummy variables.
\item
  \texttt{block\_order}: categorical, represented as 7 dummy variables.
\end{itemize}

Within each unique region, there are 6300 texture observations and 600
blank observations. Omitting \texttt{block\_order}, there are 90
repeated observations for each combination of \texttt{coherence} and
\texttt{family}. The distribution of \texttt{block\_order} is not
uniform for these combinations.

\subsection{Potential predictors}\label{potential-predictors}

Unfortunately, according to Gerick, the parameters used in the texture
generation model are probably lost. However, he mentioned that it should
be possible to get the original stimulus images, i.e., the images
presented during the passive fixation recordings. If I get these images,
I should be able to design a few continuous predictors based off of
similar statistics to the Portilla and Simoncelli parametric texture
model, e.g., spatial correlations, luminance variance, etc.

\section{PSTHs}\label{psths}

PSTHs taken over all stimuli and sites are plotted in
Figure~\ref{fig-psths}.

\begin{figure}

\centering{

\pandocbounded{\includegraphics[keepaspectratio]{20260625_files/figure-pdf/fig-psths-1.pdf}}

}

\caption{\label{fig-psths}Averaged MUA for blank, noisy, and
naturalistic textures for each region.}

\end{figure}%

\section{Onset times}\label{onset-times}

\section{Naive baseline model}\label{naive-baseline-model}

\subsection{Demo on single region (V1)}\label{demo-on-single-region-v1}




\end{document}
